\section{Background}

\subsection{Protocol Buffers Wire Format}
Protocol Buffers define a compact binary format in which each field is prefixed by a key encoded as a variable-length integer (varint):
\begin{equation}
\text{Key} = (\text{Field Number} \ll 3) \mid \text{Wire Type}
\label{eq:wire_key}
\end{equation}
Four wire types remain live in proto3: varint (\texttt{int32}, \texttt{int64}, \texttt{uint32}, \texttt{uint64}, \texttt{bool}, enums), 64-bit and 32-bit fixed (\texttt{fixed64}/\texttt{double} and \texttt{fixed32}/\texttt{float}), and length-delimited (\texttt{string}, \texttt{bytes}, submessages, and packed repeated scalars). Signed integers (\texttt{sint32}, \texttt{sint64}) apply Zigzag encoding, $(n \ll 1) \oplus (n \gg k-1)$ for a $k$-bit width, so that small negative values do not occupy the full varint width.

\subsection{C++26 Language Facilities}
Two ISO C++26 proposals make compile-time IDL ingestion possible. P1967 \cite{p1967} adds the \texttt{\#embed} directive, exposing an external file as a contiguous compile-time byte span without the external build scripts (\texttt{xxd -i} and similar) previously required. P2996 \cite{p2996} adds static reflection around the opaque \texttt{std::meta::info} type, supporting introspection (\texttt{\^{}\^{}T}, \texttt{members\_of()}), programmatic generation of member descriptors (\texttt{define\_aggregate()}, \texttt{data\_member\_spec()}), and splicing of reflected entities into concrete declarations (\texttt{[: r :]}). Combined with \texttt{constexpr} evaluation, the two turn the compiler front-end into a programmable meta-compiler capable of digesting an external domain-specific language and emitting native types.
